\documentclass{hw-template}

\title{NE523 Homework 1}
\author{Kyle Hansen}
\date{2 February 2026}

\usepackage{cancel}

\renewcommand{\defeq}{\stackrel{\text{def}}{=}}




\begin{document}

\maketitle{}

\section{Two-species Boltzmann Equation}

The Boltzmann equation can be derived by computing the transformation that a small element in phase space sees in a short time $\delta t$. From \autoref{sec:volume}, in the absence of collisions and considering only first-order $\delta t$ terms, the particles at $(\vec{r}, \vec{v})$ will stream to {$\left(\vec{r} + \vec{v}\delta t, \vec{v} + \frac{\vec{F}|_{t=t_0}}{m} \delta t, t_0 + \delta t\right)$} after a short time $\delta t$, or 

\begin{equation}
    f\left(\vec{r} + \vec{v}\delta t, \vec{v} + \frac{\vec{F}|_{t=t_0}}{m} \delta t, t_0 + \delta t\right)= f\left(\vec{r}, \vec{v}, t_0 \right).
\end{equation}

Note that $\delta t$ must be small enough to neglect the higher-order terms. The effect of scattering can be denoted with $\left( \frac{\partial f}{\partial t}\right)_{coll}$, the net change in $f$ over time due to scattering (into and out of $(\vec{r}, \vec{v})$). Then in time $\delta t$ the total change in $f$ is $\left( \frac{\partial f}{\partial t}\right)_{coll} \delta t$, and

\begin{equation}
    f\left(\vec{r} + \vec{v}\delta t, \vec{v} + \frac{\vec{F}|_{t=t_0}}{m} \delta t, t_0 + \delta t\right)= f\left(\vec{r}, \vec{v}, t_0 \right) + \left( \frac{\partial f}{\partial t}\right)_{coll}\delta t.
\end{equation}

Expanding the left-hand side as a Taylor series in $t$, this becomes

\begin{gather}
    \cancel{f\left(\vec{r}, \vec{v}, t \right)} + \left(\nabla_r f \cdot \vec{v} \right)\delta t + \left(\nabla_v f \cdot \frac{\vec{F}}{m} \right)\delta t + \left(\frac{\partial f}{\partial t} \right)\delta t= \cancel{f\left(\vec{r}, \vec{v}, t \right)} + \left( \frac{\partial f}{\partial t}\right)_{coll} \delta t\\
    \left(\vec{v} \nabla_r + \frac{\vec{F}}{m} \nabla_v  +\frac{\partial}{\partial t} \right)f\left(\vec{r}, \vec{v}, t \right)\delta t=  \left( \frac{\partial f}{\partial t}\right)_{coll} \delta t\\
    \left(\vec{v} \nabla_r + \frac{\vec{F}}{m} \nabla_v  +\frac{\partial}{\partial t} \right)f\left(\vec{r}, \vec{v}, t \right)=  \left( \frac{\partial f}{\partial t}\right)_{coll}.
\end{gather}

The same applies for particles of species $g$, so 

\begin{align}
    \left(\vec{v} \nabla_r + \frac{\vec{F}}{m} \nabla_v  +\frac{\partial}{\partial t} \right)f\left(\vec{r}, \vec{v}, t \right)&=  \left( \frac{\partial f}{\partial t}\right)_{coll}\\
    \left(\vec{v} \nabla_r + \frac{\vec{G}}{m} \nabla_v  +\frac{\partial}{\partial t} \right)g\left(\vec{r}, \vec{v}, t \right)&=  \left( \frac{\partial g}{\partial t}\right)_{coll}.
\end{align}

The collision integral can be found by computing the rate of scattering into or out of $(\vec{r}, \vec{v})$ due to collisions of particles of both species.



% TODO : scattering rates of two species

\section{Differential Volume Invariant}\label{sec:volume}



% TODO


\section{Maxwell-Boltzmann Moments}

The distribution with arbitrary bulk velocity $\vec{v_0}$ is

\begin{equation}
    f_0(\vec{v}) = n{\left( \frac{m}{2\pi k T} \right)}^{3/2} \exp \left[ -\frac{m |\vec{v} - \vec{v_0}|^2}{2kT}\right].
\end{equation}

Let $\vec{w} \defeq \vec{v} - \vec{v_0}$. Then distribution in the transformed coordinate is 


\begin{equation}
    f_0(\vec{w} + \vec{v_0}) = n{\left( \frac{m}{2\pi k T} \right)}^{3/2} \exp \left[ -\frac{m |\vec{w}|^2}{2kT}\right].
\end{equation}

Then the transormation dissapears when integrating the zeroth moment, since the bounds of integration are infinite: $\int d^w f_0(\vec{w} + \vec{v_0}) = \int d^3 w f_0(\vec{w})$. Then the zeroth moment is

\begin{equation}
    \int d^3w f_0(\vec{w}) = \int d^3w f_0(\vec{w} + \vec{v_0}) = \int d^3w n{\left( \frac{m}{2\pi k T} \right)}^{3/2} \exp \left[ -\frac{m |\vec{w}|^2}{2kT}\right].
\end{equation}

Transforming to spherical coordinates, the differential volume element becomes $d^3w = w^2 \sin \theta d\theta d\phi dw$, where $0\leq \theta \leq \pi$ is the polar angle, $0\leq \phi \leq 2\pi$ is the azimuthal angle, and $0 \leq (w = |\vec{w}|) < \infty$ is the speed, and the zeroth moment is

\begin{subequations}
    \begin{align}
    \int d^3w f_0(\vec{w})  &= \int_{0}^{\pi} \sin \theta d \theta \int_{0}^{2\pi} d\phi \int_{0}^{\infty}   n{\left( \frac{m}{2\pi k T} \right)}^{3/2} \exp \left[ -\frac{m w^2}{2kT}\right]w^2 dw\\
        &= 4\pi n{\left( \frac{m}{2\pi k T} \right)}^{3/2}\int_{0}^{\infty}    \exp \left[ -\frac{m w^2}{2kT}\right]w^2 dw.
\end{align}
\end{subequations}

With a change of variables $t = \frac{m w^2}{2kT}$, this becomes

\begin{equation}
    \int d^3w f_0(\vec{w}) = 4\pi n{\left( \frac{m}{2\pi k T} \right)}^{3/2}\int_{0}^{\infty} \frac{1}{2}  \left( \frac{2kT}{m} \right)^{3/2} t^{1/2}  \exp \left[ -t\right] dt.
\end{equation}

This integral can be evaluated using the gamma function $\Gamma(z) = \int_{0}^{\infty} t^{z-1}e^{-t}dt$, 

\begin{align}
    \int d^3w f_0(\vec{w}) &= 4\pi n{\left( \frac{m}{2\pi k T} \right)}^{3/2}\frac{1}{2}  \left( \frac{2kT}{m} \right)^{3/2} \Gamma\left( \frac{3}{2}\right)\\
    &= 4\pi n{\left( \frac{m}{2\pi k T} \right)}^{3/2}\frac{1}{2}  \left( \frac{2kT}{m} \right)^{3/2} \frac{\sqrt{\pi}}{2},
\end{align}

so the zeroth moment is


\begin{equation}
    \boxed{\int d^3w f_0(\vec{w}) = n}.
\end{equation}



The first moment is

\begin{multline}
    \int d^3w (\vec{w} + \vec{v_0}) f_0(\vec{w} + \vec{v_0}) = \int d^3w \left(\vec{w} + \vec{v_0} \right) f_0 (\vec{w} + \vec{v_0}) \\
    = \int d^3w \vec{w} f_0 (\vec{w} + \vec{v_0}) + \int d^3 w \vec{v_0} f_0 (\vec{w} + \vec{v_0}).
\end{multline}

The second part of this integral can be evaluated easily using the zeroth moment, since $\vec{v_0}$ does not depend on $\vec{w}$; $\int d^3 w \vec{v_0} f_0 (\vec{w} + \vec{v_0}) = \vec{v_0}n$. The first part of this integral evaluates to zero, since the bracketed term in 

\begin{multline}
    \int d^3w \vec{w} f_0 (\vec{w} + \vec{v_0})  \\
    =  n{\left( \frac{m}{2\pi k T} \right)}^{3/2}\int_{0}^{\infty}    \exp \left( -\frac{m w^2}{2kT}\right)w^3 dw \left[\int_{0}^{\pi} \sin\theta d \theta \int_{0}^{2\pi}d \phi \hat{w}(\theta, \phi)\right],
\end{multline}

where $\hat{w}(\theta, \phi)$ is the unit vector that points in the direction $(\theta, \phi)$, is equal to zero due to spherical symmetry.

Then the first moment is

\begin{equation}
    \boxed{\int d^3 v \vec{v}f_0(\vec{v}) = n\vec{v_0}}.
\end{equation}

The second moment is

\begin{equation}
    \int d^3 v \left(\frac{1}{2} mv^2 \right)    f_0(\vec{v}).
\end{equation}

Applying the same transformation $\vec{w} = \vec{v} - \vec{v_0}$, this can be written


\begin{equation}
  \int d^3 v \left(\frac{1}{2} mv^2 \right)    f_0(\vec{v})= n {\left( \frac{m}{2\pi k T} \right)}^{3/2}  \int d^3 w \left(\frac{1}{2}m |\vec{w} + \vec{v_0}|^2 \right) \exp \left[ \frac{-m w^2}{2kT} \right],
\end{equation}

where

\begin{subequations}
 \begin{align}
    |\vec{w} + \vec{v_0}|^2 &= {\left(w_x + v_{0,x} \right)}^2 + {\left(w_y + v_{0,y} \right)}^2 + {\left(w_z, v_{0, z} \right)}^2\\
    &=\left(w_x^2 + 2w_x v_{0, x} + v_{0, x}^2\right) + \left(w_y^2 + 2w_y v_{0, y} + v_{0, y}^2\right) + \left(u_z^2 + 2u_z v_{0, z} + v_{0, z}^2\right)
\end{align}   
\end{subequations}

Then, translating to spherical coordinates, this becomes

\begin{multline}
    |\vec{w} + \vec{v_0}|^2 =\left(w_x^2 + 2\sin\theta\cos\phi \cdot w v_{0, x} + v_{0, x}^2\right)\\
     + \left(w_y^2 + 2 \sin\theta\sin\phi \cdot w v_{0, y} + v_{0, y}^2\right)\\
      + \left(u_z^2 + 2 \cos\phi \cdot wv_{0, z} + v_{0, z}^2\right).
\end{multline}

Then, rewriting the integral in spherical coordinates yields

\begin{multline}
    n {\left( \frac{m}{2\pi k T} \right)}^{3/2}  \left[\int_{0}^{\pi} \sin \theta  d\theta \int_{0}^{2\pi} d \phi \int_{0}^{\infty} dw \left(\frac{1}{2}m w^4 \right) \exp \left( \frac{-m w^2}{2kT} \right) \right. \\
    +\int_{0}^{\pi}  d\theta \int_{0}^{2\pi} d \phi \int_{0}^{\infty} dw\left(m \left(\sin^2\theta \cos \phi v_{0, x} + \sin^2\theta \sin\phi v_{0, y} + \sin\theta \cos\phi v_{0, z} \right)w^3 \right) \exp \left( \frac{-m w^2}{2kT} \right)   \\
    \left.+\frac{1}{2}m v_0^2 \int_{0}^{\pi} \sin\theta d\theta \int_{0}^{2\pi} d \phi \int_{0}^{\infty}dw \left( w^2 \right) \exp \left( \frac{-m w^2}{2kT} \right)  \right].
\end{multline}

The second line of the bracketed term evaluates to zero since $\int_{0}^{2\pi} \sin \phi d \phi = \int_{0}^{2\pi}\cos \phi d \phi = 0$. With the change of variables $t = \frac{mw^2}{2kT}$, this equation becomes

\begin{subequations}
    \begin{multline}
        n {\left( \frac{m}{2\pi k T} \right)}^{3/2} \cdot 4\pi \left[\int_{0}^{\infty} \frac{1}{2}m\left(\frac{2kTt}{m} \right)^2 \exp(-t) \cdot \frac{1}{2} \left( \frac{2kT}{m}\right)^{1/2} t^{1/2}dt \right. \\
        + \frac{1}{2}m v_0^2\int_{0}^{\infty} \frac{1}{2}m\left(\frac{2kTt}{m} \right) \exp(-t) \cdot \frac{1}{2} \left( \frac{2kT}{m}\right)^{1/2} t^{1/2}dt
\end{multline}

%TODO-- THIS INTEGRAL IS WRONG

\end{subequations}



% TODO  


\end{document}
